\documentclass[10pt,a4paper]{amsart}
\usepackage{amsmath,amssymb,amsthm,graphicx}
\usepackage{verbatim}
\newenvironment{reference}{\comment}{\endcomment}
%\newenvironment{reference}{}{}
\newenvironment{slogan}{\comment}{\endcomment}
\newenvironment{history}{\comment}{\endcomment}

% For commutative diagrams we use Xy-pic
\usepackage[all]{xy}

% We use 2cell for 2-commutative diagrams.
\xyoption{2cell}
% Theorem environments.
%
\theoremstyle{plain}
\newtheorem{theorem}[subsection]{Theorem}
\newtheorem{proposition}[subsection]{Proposition}
\newtheorem{lemma}[subsection]{Lemma}

\theoremstyle{definition}
\newtheorem{definition}[subsection]{Definition}
\newtheorem{example}[subsection]{Example}
\newtheorem{exercise}[subsection]{Exercise}
\newtheorem{situation}[subsection]{Situation}

\theoremstyle{remark}
\newtheorem{remark}[subsection]{Remark}
\newtheorem{remarks}[subsection]{Remarks}

\numberwithin{equation}{subsection}

% Macros
%
\def\lim{\mathop{\mathrm{lim}}\nolimits}
\def\colim{\mathop{\mathrm{colim}}\nolimits}
\def\Spec{\mathop{\mathrm{Spec}}}
\def\Hom{\mathop{\mathrm{Hom}}\nolimits}
\def\Ext{\mathop{\mathrm{Ext}}\nolimits}
\def\SheafHom{\mathop{\mathcal{H}\!\mathit{om}}\nolimits}
\def\SheafExt{\mathop{\mathcal{E}\!\mathit{xt}}\nolimits}
\def\Sch{\mathit{Sch}}
\def\Mor{\mathop{\mathrm{Mor}}\nolimits}
\def\Ob{\mathop{\mathrm{Ob}}\nolimits}
\def\Sh{\mathop{\mathit{Sh}}\nolimits}
\def\NL{\mathop{N\!L}\nolimits}
\def\CH{\mathop{\mathrm{CH}}\nolimits}
\def\proetale{{pro\text{-}\acute{e}tale}}
\def\etale{{\acute{e}tale}}
\def\QCoh{\mathit{QCoh}}
\def\Ker{\mathop{\mathrm{Ker}}}
\def\Im{\mathop{\mathrm{Im}}}
\def\Coker{\mathop{\mathrm{Coker}}}
\def\Coim{\mathop{\mathrm{Coim}}}

% Boxtimes
%
\DeclareMathSymbol{\boxtimes}{\mathbin}{AMSa}{"02}

%
% Macros for moduli stacks/spaces
%
\def\QCohstack{\mathcal{QC}\!\mathit{oh}}
\def\Cohstack{\mathcal{C}\!\mathit{oh}}
\def\Spacesstack{\mathcal{S}\!\mathit{paces}}
\def\Quotfunctor{\mathrm{Quot}}
\def\Hilbfunctor{\mathrm{Hilb}}
\def\Curvesstack{\mathcal{C}\!\mathit{urves}}
\def\Polarizedstack{\mathcal{P}\!\mathit{olarized}}
\def\Complexesstack{\mathcal{C}\!\mathit{omplexes}}
% \Pic is the operator that assigns to X its picard group, usage \Pic(X)
% \Picardstack_{X/B} denotes the Picard stack of X over B
% \Picardfunctor_{X/B} denotes the Picard functor of X over B
\def\Pic{\mathop{\mathrm{Pic}}\nolimits}
\def\Picardstack{\mathcal{P}\!\mathit{ic}}
\def\Picardfunctor{\mathrm{Pic}}
\def\Deformationcategory{\mathcal{D}\!\mathit{ef}}
\usepackage{cleveref}
\crefname{originallemma}{original lemma}{original lemmas}
\Crefname{originallemma}{Original lemma}{Original lemmas}
\makeatletter
\newlabel{algebra-lemma-length-independent}{{00IX}{source}{}{sourcepointer-1}{}}
\newlabel{algebra-lemma-length-independent@cref}{{[originallemma][0][]00IX}{[1][0][]source}{}{}{}}
\makeatother
\crefname{originalequation}{original equation}{original equations}
\Crefname{originalequation}{Original equation}{Original equations}
\makeatletter
\newlabel{spaces-cohomology-equation-representable-higher-direct-image}{{071Z}{source}{}{sourcepointer-2}{}}
\newlabel{spaces-cohomology-equation-representable-higher-direct-image@cref}{{[originalequation][0][]071Z}{[1][0][]source}{}{}{}}
\makeatother
\crefname{originallemma}{original lemma}{original lemmas}
\Crefname{originallemma}{Original lemma}{Original lemmas}
\makeatletter
\newlabel{spaces-lemma-base-change-immersions}{{02YW}{source}{}{sourcepointer-3}{}}
\newlabel{spaces-lemma-base-change-immersions@cref}{{[originallemma][0][]02YW}{[1][0][]source}{}{}{}}
\makeatother
\crefname{originallemma}{original lemma}{original lemmas}
\Crefname{originallemma}{Original lemma}{Original lemmas}
\makeatletter
\newlabel{spaces-properties-lemma-points-cartesian}{{03H4}{source}{}{sourcepointer-4}{}}
\newlabel{spaces-properties-lemma-points-cartesian@cref}{{[originallemma][0][]03H4}{[1][0][]source}{}{}{}}
\makeatother
\title{The length of a quasi-coherent module at a point is unchanged by pushforward along a closed immersion}
\author{The Stacks Project Authors}
\date{stacks-project:spaces-chow.tex:local-snapshot-20261012; GNU FDL 1.2 or later}
\begin{document}
\maketitle
\noindent\textit{Editorial note.} All delivered mathematical text is the author's own, copied from the pinned Stacks Project chapter file spaces-chow.tex; no statement or proof step was rewritten, reordered or removed.\par\medskip
\noindent\textit{Editorial note.} The counted proof is the author's complete argument: an etale chart reduction to closed immersions of schemes, an identification of the restricted pushforward with the pushforward of the restricted module, and a length comparison through the surjection of local rings\par\medskip
\noindent\textit{Editorial note.} The label lemma-length-closed-immersion names the counted statement, and the proof environment immediately following it in the source is delivered verbatim as its proof.\par\medskip
\noindent\textit{Editorial note.} The counted statement is carried by the author's own lemma environment in the chapter source.\par\medskip
\noindent\textit{Allowed standard prerequisites (editorial).}\begin{itemize}
\item The length of a quasi-coherent module at a point, and its compatibility with etale charts
\item Etale charts and the cartesian description of the points of a fibre product, together with stability of closed immersions under base change
\item The identification of the pushforward along a base changed closed immersion with the restriction of the pushforward, and the algebra statement that the length of a module over a local ring agrees with its length viewed through a surjection of local rings
\end{itemize}\medskip
\ifdefined\hypertarget\hypertarget{sourcepointer-1}{}\fi
\noindent\textit{Source reference.} Original lemma 00IX: sources/stacks-chow/tags.txt, source lines 695--695. Prerequisite pointer invoked at proof line(s) 555; it names a standard fact of a sibling chapter that the counted argument uses, and every step of the counted proof built on it is retained verbatim in the delivered body.\par\medskip
\ifdefined\hypertarget\hypertarget{sourcepointer-2}{}\fi
\noindent\textit{Source reference.} Original equation 071Z: sources/stacks-chow/tags.txt, source lines 9008--9008. Prerequisite pointer invoked at proof line(s) 545; it names a standard fact of a sibling chapter that the counted argument uses, and every step of the counted proof built on it is retained verbatim in the delivered body.\par\medskip
\ifdefined\hypertarget\hypertarget{sourcepointer-3}{}\fi
\noindent\textit{Source reference.} Original lemma 02YW: sources/stacks-chow/tags.txt, source lines 3891--3891. Prerequisite pointer invoked at proof line(s) 540; it names a standard fact of a sibling chapter that the counted argument uses, and every step of the counted proof built on it is retained verbatim in the delivered body.\par\medskip
\ifdefined\hypertarget\hypertarget{sourcepointer-4}{}\fi
\noindent\textit{Source reference.} Original lemma 03H4: sources/stacks-chow/tags.txt, source lines 4553--4553. Prerequisite pointer invoked at proof line(s) 539; it names a standard fact of a sibling chapter that the counted argument uses, and every step of the counted proof built on it is retained verbatim in the delivered body.\par\medskip
\begin{lemma}
\label{lemma-length-closed-immersion}
Let $S$ be a scheme. Let $i : Y \to X$ be a closed immersion of
algebraic spaces over $S$. Let $\mathcal{G}$ be a quasi-coherent
$\mathcal{O}_Y$-module. Let $y \in |Y|$ with image $x \in |X|$.
Let $d \in \{0, 1, 2, \ldots, \infty\}$. The following are
equivalent
\begin{enumerate}
\item $\mathcal{G}$ has length $d$ at $y$, and
\item $i_*\mathcal{G}$ has length $d$ at $x$.
\end{enumerate}
\end{lemma}

\begin{proof}
Choose an \'etale morphism $f : U \to X$ with $U$ a scheme
and $u \in U$ mapping to $x$. Set $V = Y \times_X U$.
Denote $g : V \to Y$ and $j : V \to U$ the projections.
Then $j : V \to U$ is a closed immersion and there is a unique
point $v \in V$ mapping to $y \in |Y|$ and $u \in U$
(use Properties of Spaces, Lemma \ref{spaces-properties-lemma-points-cartesian}
and Spaces, Lemma \ref{spaces-lemma-base-change-immersions}).
We have $j_*(\mathcal{G}|_V) = (i_*\mathcal{G})|_U$ as modules
on the scheme $V$ and $j_*$ the ``usual'' pushforward of modules
for the morphism of schemes $j$, see discussion surrounding
Cohomology of Spaces, Equation
(\ref{spaces-cohomology-equation-representable-higher-direct-image}).
In this way we reduce to the case of schemes: if $i : Y \to X$
is a closed immersion of schemes, then
$$
(i_*\mathcal{G})_x = \mathcal{G}_y
$$
as modules over $\mathcal{O}_{X, x}$ where the module structure
on the right hand side is given by the surjection
$i_y^\sharp : \mathcal{O}_{X, x} \to \mathcal{O}_{Y, y}$.
Thus equality by
Algebra, Lemma \ref{algebra-lemma-length-independent}.
\end{proof}

\end{document}
